\chapter{2013年浙江卷（理）}

\section{单选题}

\begin{problem}
已知 \(\i\) 是虚数单位，则 \((-1 + \i)(2 - \i) =\)（\quad）

\choices
    {\(-3 + \i\)}
    {\(-1 + 3\i\)}
    {\(-3 + 3\i\)}
    {\(-1 + \i\)}
\end{problem}

\begin{answer}
B
\end{answer}

\begin{solution}
\( (-1+\i)(2-\i)=-2+\i+2\i-\i^{2}=-1+3\i\)，故选 B.
\end{solution}

\begin{problem}
设集合 \(S = \{x \mid x > -2\}\)，\(T = \{x \mid x^2 + 3x - 4 \leqslant 0\}\)，则 \(\left(\complement_{\R} S\right) \cup T =\)（\quad）

\choices
    {\((-2, 1]\)}
    {\((-\infty, -4]\)}
    {\((-\infty, 1]\)}
    {\([1, +\infty)\)}
\end{problem}

\begin{answer}
C
\end{answer}

\begin{solution}
由 \(x^{2}+3x-4=(x+4)(x-1)\)，得 \(T=[-4,1]\)，又 \(\complement_{\R}S=(-\infty,-2]\)，所以 \(\left(\complement_{\R}S\right)\cup T=(-\infty,1]\)，故选 C.
\end{solution}

\begin{problem}
已知 \(x\)，\(y\) 为正实数，则（\quad）

\choices
    {\(2^{\lg x + \lg y} = 2^{\lg x} + 2^{\lg y}\)}
    {\(2^{\lg (x + y)} = 2^{\lg x} \cdot 2^{\lg y}\)}
    {\(2^{\lg x - \lg y} = 2^{\lg x} + 2^{\lg y}\)}
    {\(2^{\lg (xy)} = 2^{\lg x} \cdot 2^{\lg y}\)}
\end{problem}

\begin{answer}
D
\end{answer}

\begin{solution}
因为 \(x\)，\(y\) 为正实数，所以 \(\lg(xy)=\lg x+\lg y\)，从而 \(2^{\lg(xy)}=2^{\lg x+\lg y}=2^{\lg x}\cdot2^{\lg y}\)，故选 D.
\end{solution}

\begin{problem}
已知函数 \( f(x) = A\cos (\omega x + \varphi)  \) （\(A > 0\)，\(\omega > 0\)，\(\varphi \in \R\)），则“\(f(x)\) 是奇函数”是“\( \varphi = \frac{\pi}{2} \)”的（\quad）

\choices
    {充分不必要条件}
    {必要不充分条件}
    {充分必要条件}
    {既不充分也不必要条件}
\end{problem}

\begin{answer}
B
\end{answer}

\begin{solution}
若 \(f(x)\) 是奇函数，则 \(f(0)=A\cos\varphi=0\)，所以 \(\varphi=\dfrac{\pi}{2}+k\pi\)，其中 \(k\in\Z\). 当 \(\varphi=\dfrac{\pi}{2}\) 时，\(f(x)=-A\sin\omega x\) 是奇函数. 因此，“\(f(x)\) 是奇函数”是“\(\varphi=\dfrac{\pi}{2}\)”的必要不充分条件，故选 B.
\end{solution}

\begin{problem}
某程序框图如图所示，若该程序运行后输出的值是 \(\frac{9}{5}\)，则（\quad）

\texfigure[max width=7.0cm]{img_repaint/2013/zhejiang_science/q5_fig1.tex}

\choices
    {\(a = 4\)}
    {\(a = 5\)}
    {\(a = 6\)}
    {\(a = 7\)}
\end{problem}

\begin{answer}
A
\end{answer}

\begin{solution}
程序循环结束时，循环变量 \(k\) 已从 \(1\) 依次取到 \(a\)，因此输出值为
\(S=1+\displaystyle\sum_{k=1}^{a}\dfrac{1}{k(k+1)}=1+\displaystyle\sum_{k=1}^{a}\left(\dfrac{1}{k}-\dfrac{1}{k+1}\right)=2-\dfrac{1}{a+1}\). 由输出值为 \(\dfrac{9}{5}\)，得 \(2-\dfrac{1}{a+1}=\dfrac{9}{5}\)，即 \(a+1=5\)，所以 \(a=4\)，故选 A.
\end{solution}

\begin{problem}
已知 \(\alpha \in \R\)，\(\sin \alpha + 2\cos \alpha = \frac{\sqrt{10}}{2}\)，则 \(\tan 2\alpha =\)（\quad）

\choices
    {\(\frac{4}{3}\)}
    {\( \frac{3}{4} \)}
    {\(-\frac{3}{4}\)}
    {\(-\frac{4}{3}\)}
\end{problem}

\begin{answer}
C
\end{answer}

\begin{solution}
由已知式得 \(\cos\alpha\ne0\). 令 \(t=\tan\alpha\)，则
\(\sin\alpha+2\cos\alpha=(t+2)\cos\alpha=\dfrac{\sqrt{10}}{2}\). 两边平方并利用 \(\cos^{2}\alpha=\dfrac{1}{1+t^{2}}\)，得 \(\dfrac{(t+2)^{2}}{1+t^{2}}=\dfrac{5}{2}\)，因此 \(3t^{2}-8t-3=0\)，所以 \(t=3\) 或 \(t=-\dfrac{1}{3}\). 对于 \(t=3\)，原式要求 \(\cos\alpha=\dfrac{1}{\sqrt{10}}\)；对于 \(t=-\dfrac{1}{3}\)，原式要求 \(\cos\alpha=\dfrac{3}{\sqrt{10}}\)，两种情形均可满足原式，故平方没有引入不合题意的根. 于是两种情形下均有 \(\tan2\alpha=\dfrac{2t}{1-t^{2}}=-\dfrac{3}{4}\)，故选 C.
\end{solution}

\begin{problem}
设 \(\triangle ABC\)，\(P_0\) 是边 \(AB\) 上一定点. 满足 \(P_0B = \frac{1}{4}AB\)，且对于边 \(AB\) 上任一点 \(P\)，恒有 \(\overrightarrow{PB} \cdot \overrightarrow{PC} \geqslant \overrightarrow{P_0B} \cdot \overrightarrow{P_0C}\). 则（\quad）

\choices
    {\(\angle ABC = 90^{\circ}\)}
    {\(\angle BAC = 90^{\circ}\)}
    {\(AB = AC\)}
    {\(AC = BC\)}
\end{problem}

\begin{answer}
D
\end{answer}

\begin{solution}
设 \(AB=L\)，取 \(B=(0,0)\)，\(A=(L,0)\)，并设 \(C=(u,v)\). 边 \(AB\) 上的点可写为 \(P=(t,0)\)，其中 \(0\leqslant t\leqslant L\). 于是
\(\overrightarrow{PB}\cdot\overrightarrow{PC}=(-t,0)\cdot(u-t,v)=t(t-u)\). 由 \(P_{0}B=\dfrac{L}{4}\)，点 \(P_{0}\) 对应 \(t_{0}=\dfrac{L}{4}\)，且 \(t_{0}\) 是区间内部点. 二次函数 \(t(t-u)\) 在 \(t_{0}\) 处取得区间上的最小值，所以其导数在此处为零，即 \(2t_{0}-u=0\)，从而 \(u=\dfrac{L}{2}\). 因此 \(BC^{2}=u^{2}+v^{2}=(u-L)^{2}+v^{2}=AC^{2}\)，故 \(AC=BC\)，选 D.
\end{solution}

\begin{problem}
已知 \(\e\) 为自然对数的底数，设函数 \( f(x) = (\e^{x} - 1)(x - 1)^{k}  \) （\(k = 1\)，\(2\)），则（\quad）

\choices
    {当 \(k = 1\) 时，\(f(x)\) 在 \(x = 1\) 处取得极小值}
    {当 \(k = 1\) 时，\(f(x)\) 在 \(x = 1\) 处取得极大值}
    {当 \(k = 2\) 时，\(f(x)\) 在 \(x = 1\) 处取得极小值}
    {当 \(k = 2\) 时，\(f(x)\) 在 \(x = 1\) 处取得极大值}
\end{problem}

\begin{answer}
C
\end{answer}

\begin{solution}
在 \(x=1\) 附近，\(\e^{x}-1>0\). 当 \(k=1\) 时，\(f(x)=(\e^{x}-1)(x-1)\) 在 \(x=1\) 的两侧符号相反，故在此处既不取得极大值也不取得极小值. 当 \(k=2\) 时，在 \(x=1\) 的某邻域内有 \(f(x)=(\e^{x}-1)(x-1)^{2}\geqslant0\)，且 \(f(1)=0\)，所以在 \(x=1\) 处取得极小值，故选 C.
\end{solution}

\begin{problem}
如图，\(F_{1}\)，\(F_{2}\) 是椭圆 \(C_{1}: \frac{x^{2}}{4} + y^{2} = 1\) 与双曲线 \(C_{2}\) 的公共焦点，\(A\)，\(B\) 分别是 \(C_{1}\)，\(C_{2}\) 在第二，四象限的公共点. 若四边形 \(AF_{1}BF_{2}\) 为矩形. 则 \(C_{2}\) 的离心率是（\quad）

\FigureLayoutDeclare{img/2013/zhejiang_science/q9_fig1.png}{6.6cm}{22bp 19bp 8bp 20bp}
\bitmapfigure[max width=6.0cm]{img/2013/zhejiang_science/q9_fig1.png}

\choices
    {\(\sqrt{2}\)}
    {\(\sqrt{3}\)}
    {\( \frac{3}{2} \)}
    {\(\frac{\sqrt{6}}{2}\)}
\end{problem}

\begin{answer}
D
\end{answer}

\begin{solution}
椭圆 \(C_{1}\) 的半长轴、半短轴分别为 \(2\)，\(1\)，故其焦点为 \(F_{1}(-\sqrt{3},0)\)，\(F_{2}(\sqrt{3},0)\). 设 \(A=(u,v)\)，则由矩形的两条对角线互相平分，得 \(B=(-u,-v)\). 又相邻两边垂直，所以
\(\overrightarrow{AF_{1}}\cdot\overrightarrow{F_{1}B}=(-\sqrt{3}-u,-v)\cdot(\sqrt{3}-u,-v)=u^{2}+v^{2}-3=0\). 结合 \(A\) 在椭圆上，得 \(u^{2}+v^{2}=3\)，\(u^{2}+4v^{2}=4\)，从而 \(v^{2}=\dfrac{1}{3}\)，\(u^{2}=\dfrac{8}{3}\).

设双曲线 \(C_{2}\) 的方程为 \(\dfrac{x^{2}}{p^{2}}-\dfrac{y^{2}}{q^{2}}=1\)，则 \(p^{2}+q^{2}=3\). 点 \(A\) 在双曲线上，故 \(\dfrac{8}{3p^{2}}-\dfrac{1}{3q^{2}}=1\). 令 \(p^{2}=m\)，则 \(q^{2}=3-m\)，代入得 \(\dfrac{8}{m}-\dfrac{1}{3-m}=3\)，解得 \(m=2\) 或 \(m=4\). 因 \(q^{2}>0\)，只能取 \(p^{2}=2\)，于是 \(C_{2}\) 的离心率为 \(e=\dfrac{\sqrt{p^{2}+q^{2}}}{p}=\dfrac{\sqrt{3}}{\sqrt{2}}=\dfrac{\sqrt{6}}{2}\)，故选 D.
\end{solution}

\begin{problem}
在空间中，过点 \(A\) 作平面 \(\pi\) 的垂线，垂足为 \(B\). 记 \(B = f_{\pi}(A)\). 设 \(\alpha\)，\(\beta\) 是两个不同的平面，对空间任意一点 \(P\)，\(Q_1 = f_{\beta}(f_{\alpha}(P))\)，\(Q_2 = f_{\alpha}(f_{\beta}(P))\)，恒有 \(PQ_1 = PQ_2\)，则（\quad）

\choices
    {平面 \(\alpha\) 与平面 \(\beta\) 垂直}
    {平面 \(\alpha\) 与平面 \(\beta\) 所成的（锐）二面角为 \(45^{\circ}\)}
    {平面 \(\alpha\) 与平面 \(\beta\) 平行}
    {平面 \(\alpha\) 与平面 \(\beta\) 所成的（锐）二面角为 \(60^{\circ}\)}
\end{problem}

\begin{answer}
A
\end{answer}

\begin{solution}
若 \(\alpha\)，\(\beta\) 平行且不同，取 \(P\in\alpha\)，则 \(Q_{2}=f_{\alpha}(f_{\beta}(P))=P\)，而 \(Q_{1}\) 是 \(P\) 在 \(\beta\) 上的垂足，故 \(PQ_{1}>0=PQ_{2}\)，与已知矛盾. 因此两平面相交.

在垂直于两平面交线的截面内，设两平面的截线分别为两条相交直线，取第一条为 \(x\) 轴，第二条的单位方向向量为 \((\cos\theta,\sin\theta)\). 任取截面内一点 \(P=(x,y)\)，记 \(c=\cos\theta\)，\(s=\sin\theta\). 两次投影所得点分别为 \(Q_{1}=(xc^{2},xcs)\) 和 \(Q_{2}=(xc^{2}+ysc,0)\). 直接计算得
\(PQ_{1}^{2}-PQ_{2}^{2}=sc^{2}\bigl(s(x^{2}-y^{2})-2cxy\bigr)\). 已知等式对任意 \(P\) 成立，取 \((x,y)=(1,0)\)，得到 \(s^{2}c^{2}=0\). 两平面相交且不同，故 \(s\ne0\)，从而 \(c=0\)，即 \(\theta=90^{\circ}\). 所以两平面垂直，选 A.
\end{solution}

\section{填空题}

\begin{problem}
设二项式 \(\left(\sqrt{x} - \frac{1}{\sqrt[3]{x}}\right)^5\) 的展开式中常数项为 \(A\)，则 \(A =\)\fillinblank{}.
\end{problem}

\begin{answer}
\(-10\)
\end{answer}

\begin{solution}
展开式的通项为 \((-1)^{r}\mathrm C_{5}^{r}x^{\frac{5-r}{2}-\frac{r}{3}}=(-1)^{r}\mathrm C_{5}^{r}x^{\frac{15-5r}{6}}\)，其中 \(r=0\)，\(1\)，\(\ldots\)，\(5\). 常数项要求 \(\dfrac{15-5r}{6}=0\)，所以 \(r=3\)，从而 \(A=(-1)^{3}\mathrm C_{5}^{3}=-10\).
\end{solution}

\begin{problem}
若某几何体的三视图（单位：\(\mathrm{cm}\)）如图所示，则此几何体的体积等于\fillinblank{}\(\mathrm{cm}^{3}\).

\bitmapfigure[max width=4.6cm]{img/2013/zhejiang_science/q12_fig1.png}
\end{problem}

\begin{answer}
24
\end{answer}

\begin{solution}
由三视图可知，该几何体可看作底面直角三角形两直角边分别为 \(3\)，\(4\)、高为 \(5\) 的直三棱柱，截去一个底面直角三角形两直角边为 \(3\)，\(4\)、高为 \(3\) 的三棱锥. 因此其体积为
\(V=\dfrac{1}{2}\times3\times4\times5-\dfrac{1}{3}\times\left(\dfrac{1}{2}\times3\times4\right)\times3=30-6=24\)，故填 \(24\).
\end{solution}

\begin{problem}
设 \(z = kx + y\)，其中实数 \(x\)，\(y\) 满足 \(\begin{cases}
x + y - 2 \geqslant 0,\\
x - 2y + 4 \geqslant 0,\\
2x - y - 4 \leqslant 0.
\end{cases}\) 若 \(z\) 的最大值为 12，则实数 \(k =\fillinblank{}\).
\end{problem}

\begin{answer}
2
\end{answer}

\begin{solution}
三个边界直线分别为 \(x+y=2\)，\(x-2y=-4\)，\(2x-y=4\). 它们两两相交所得可行域顶点为 \(A(0,2)\)，\(B(2,0)\)，\(C(4,4)\). 在三个顶点处，\(z=kx+y\) 的值分别为 \(2\)，\(2k\)，\(4k+4\).

若最大值为 \(12\)，则不可能由 \(2\) 取得；若由 \(2k\) 取得，则 \(k=6\)，但此时在 \(C\) 点有 \(z=28>12\)，矛盾. 因此只能有 \(4k+4=12\)，解得 \(k=2\).
\end{solution}

\begin{problem}
将 \(A\)，\(B\)，\(C\)，\(D\)，\(E\)，\(F\) 六个字母排成一排，且 \(A\)，\(B\) 均在 \(C\) 的同侧，则不同的排法共有\fillinblank{} 种. （用数字作答）
\end{problem}

\begin{answer}
480
\end{answer}

\begin{solution}
六个字母的全排列共有 \(6!\) 种. 固定 \(A\)，\(B\)，\(C\) 的相对次序，这三个字母的相对次序共有 \(3!=6\) 种. 当 \(C\) 在 \(A\)，\(B\) 两者左侧时，\(A\)，\(B\) 的先后次序有 \(2\) 种；当 \(C\) 在两者右侧时，\(A\)，\(B\) 的先后次序也有 \(2\) 种，所以满足条件的相对次序共有 \(4\) 种. 因此所求排法数为 \(6!\times\dfrac{4}{6}=480\).
\end{solution}

\begin{problem}
设 \(F\) 为抛物线 \(C: y^{2} = 4x\) 的焦点，过点 \(P(-1,0)\) 的直线 \(l\) 交抛物线 \(C\) 于 \(A\)，\(B\) 两点，点 \(Q\) 为线段 \(AB\) 的中点. 若 \(|FQ| = 2\)，则直线 \(l\) 的斜率等于\fillinblank{}.
\end{problem}

\begin{answer}
\(\pm 1\)
\end{answer}

\begin{solution}
若 \(l\) 为竖直线，则其方程为 \(x=-1\)，与抛物线 \(y^{2}=4x\) 没有公共点. 因而可设直线 \(l\) 的斜率为 \(m\)，且 \(l:y=m(x+1)\). 当 \(m=0\) 时，\(l\) 与抛物线只有一个交点，故 \(m\ne0\).

将 \(l\) 与抛物线联立，得 \(m^{2}(x+1)^{2}=4x\)，即 \(m^{2}x^{2}+(2m^{2}-4)x+m^{2}=0\). 设两交点的横坐标为 \(x_{1}\)，\(x_{2}\)，由韦达定理得 \(x_{1}+x_{2}=\dfrac{4}{m^{2}}-2\). 故弦中点 \(Q\) 的坐标为
\(Q\left(\dfrac{x_{1}+x_{2}}{2},\dfrac{m(x_{1}+1)+m(x_{2}+1)}{2}\right)=\left(\dfrac{2}{m^{2}}-1,\dfrac{2}{m}\right)\). 抛物线焦点为 \(F(1,0)\)，所以 \(|FQ|^{2}=\left(\dfrac{2}{m^{2}}-2\right)^{2}+\left(\dfrac{2}{m}\right)^{2}=4\left(\dfrac{1}{m^{4}}-\dfrac{1}{m^{2}}+1\right)\). 由 \(|FQ|=2\)，得 \(\dfrac{1}{m^{4}}-\dfrac{1}{m^{2}}=0\). 因 \(m\ne0\)，故 \(m^{2}=1\)，即 \(m=\pm1\).
\end{solution}

\begin{problem}
\(\triangle ABC\) 中，\(\angle C = 90^{\circ}\)，\(M\) 是 \(BC\) 的中点. 若 \(\sin \angle BAM = \frac{1}{3}\)，则 \(\sin \angle BAC =\fillinblank{}\).
\end{problem}

\begin{answer}
\(\frac{\sqrt{6}}{3}\)
\end{answer}

\begin{solution}
设 \(C=(0,0)\)，\(B=(2b,0)\)，\(A=(0,a)\)，其中 \(a>0\)，\(b>0\)，则 \(M=(b,0)\). 由向量的叉积公式，\(\sin\angle BAM=\dfrac{|\overrightarrow{AB}\times\overrightarrow{AM}|}{|AB|\,|AM|}=\dfrac{ab}{\sqrt{4b^{2}+a^{2}}\sqrt{b^{2}+a^{2}}}=\dfrac{1}{3}\). 令 \(t=\dfrac{a^{2}}{b^{2}}\)，平方后得 \(\dfrac{t}{(4+t)(1+t)}=\dfrac{1}{9}\)，即 \(t^{2}-4t+4=0\)，所以 \(t=2\).

又 \(\sin\angle BAC=\dfrac{2b}{\sqrt{4b^{2}+a^{2}}}=\dfrac{2}{\sqrt{4+t}}=\dfrac{2}{\sqrt6}=\dfrac{\sqrt6}{3}\).
\end{solution}

\begin{problem}
设 \(\bs{e}_{1}\)，\(\bs{e}_{2}\) 为单位向量，非零向量 \(\bs{b} = x\bs{e}_{1} + y\bs{e}_{2}\)，\(x\)，\(y \in \R\). 若 \(\bs{e}_{1}\)，\(\bs{e}_{2}\) 的夹角为 \(\frac{\pi}{6}\)，则 \(\frac{|x|}{|\bs{b}|}\) 的最大值等于\fillinblank{}.
\end{problem}

\begin{answer}
2
\end{answer}

\begin{solution}
由 \(\bs{e}_{1}\)，\(\bs{e}_{2}\) 为单位向量且夹角为 \(\dfrac{\pi}{6}\)，得 \(\bs{e}_{1}\cdot\bs{e}_{2}=\dfrac{\sqrt3}{2}\). 因此 \(|\bs{b}|^{2}=|x\bs{e}_{1}+y\bs{e}_{2}|^{2}=x^{2}+y^{2}+\sqrt3xy\). 若 \(x=0\)，则比值为 \(0\)；当 \(x\ne0\) 时令 \(u=\dfrac{y}{x}\)，则 \(\left(\dfrac{|x|}{|\bs{b}|}\right)^{2}=\dfrac{1}{u^{2}+\sqrt3u+1}=\dfrac{1}{\left(u+\dfrac{\sqrt3}{2}\right)^{2}+\dfrac14}\leqslant4\). 当 \(u=-\dfrac{\sqrt3}{2}\) 时等号成立，故 \(\dfrac{|x|}{|\bs{b}|}\) 的最大值为 \(2\).
\end{solution}

\section{解答题}

\begin{problem}
在公差为 \(d\) 的等差数列 \(\{a_{n}\}\) 中，已知 \(a_{1} = 10\)，且 \(a_{1}\)，\(2a_{2} + 2\)，\(5a_{3}\) 成等比数列.

\begin{enumerate}
\item 求 \(d\)，\(a_{n}\)；
\item 若 \(d < 0\)，求 \(\left|a_{1}\right| + \left|a_{2}\right| + \left|a_{3}\right| + \cdots + \left|a_{n}\right|\).
\end{enumerate}
\end{problem}

\begin{solution}
\begin{enumerate}
\item 由等差数列的通项公式，得 \(a_{2}=10+d\)，\(a_{3}=10+2d\). 由 \(a_{1}\)，\(2a_{2}+2\)，\(5a_{3}\) 成等比数列，得
\((2a_{2}+2)^{2}=a_{1}\cdot5a_{3}\)，即 \((22+2d)^{2}=10(50+10d)\). 化简得 \(d^{2}-3d-4=0\)，所以 \(d=4\) 或 \(d=-1\). 因此 \(a_{n}=10+(n-1)d\)，即 \(d=4\) 时 \(a_{n}=4n+6\)，\(d=-1\) 时 \(a_{n}=11-n\).

\item 当 \(d<0\) 时，\(d=-1\)，所以 \(a_{k}=11-k\). 当 \(1\leqslant n\leqslant11\) 时，\(a_{1}\)，\(\ldots\)，\(a_{n}\) 非负，故
\(\displaystyle\sum_{k=1}^{n}|a_{k}|=\sum_{k=1}^{n}(11-k)=\dfrac{n(21-n)}{2}\). 当 \(n\geqslant12\) 时，前 \(11\) 项绝对值之和为 \(10+9+\cdots+0=55\)，其余各项绝对值为 \(1\)，\(2\)，\(\ldots\)，\(n-11\)，所以
\[
\sum_{k=1}^{n}|a_{k}|=55+\frac{(n-11)(n-10)}{2}
\]
\end{enumerate}
\end{solution}

\begin{problem}
设袋子中装有 \(a\) 个红球，\(b\) 个黄球，\(c\) 个蓝球，且规定：取出一个红球得 \(1\) 分，取出一个黄球得 \(2\) 分，取出一个蓝球得 \(3\) 分.

\begin{enumerate}
\item 当 \(a = 3\)，\(b = 2\)，\(c = 1\) 时，从该袋子中任取（有放回，且每球取到的机会均等）2个球，记随机变量 \(\xi\) 为取出此 \(2\) 个球所得分数之和. 求 \(\xi\) 分布列；
\item 从该袋子中任取（且每球取到的机会均等）1个球，记随机变量 \(\eta\) 为取出此球所得分数. 若 \(E(\eta) = \frac{5}{3}\)，\(D(\eta) = \frac{5}{9}\)，求 \(a:b:c\).
\end{enumerate}
\end{problem}

\begin{solution}
\begin{enumerate}
\item 此时取到红球、黄球、蓝球的概率分别为 \(\dfrac12\)，\(\dfrac13\)，\(\dfrac16\)，且两次取球相互独立. 故 \(\xi\) 的可能取值为 \(2\)，\(3\)，\(4\)，\(5\)，\(6\)，相应概率分别为
\(P(\xi=2)=\left(\dfrac12\right)^{2}=\dfrac14\)，\(P(\xi=3)=2\times\dfrac12\times\dfrac13=\dfrac13\)，\(P(\xi=4)=2\times\dfrac12\times\dfrac16+\left(\dfrac13\right)^{2}=\dfrac{5}{18}\)，\(P(\xi=5)=2\times\dfrac13\times\dfrac16=\dfrac19\)，\(P(\xi=6)=\left(\dfrac16\right)^{2}=\dfrac1{36}\). 因此分布列为
\begin{center}
\begin{tabular}{c|ccccc}
\(\xi\) & \(2\) & \(3\) & \(4\) & \(5\) & \(6\) \\
\hline
\(P\) & \(\dfrac14\) & \(\dfrac13\) & \(\dfrac{5}{18}\) & \(\dfrac19\) & \(\dfrac1{36}\) \\
\end{tabular}
\end{center}

\item 设袋中球的总数为 \(N=a+b+c\). 由 \(E(\eta)=\dfrac53\)，得
\(\dfrac{a+2b+3c}{N}=\dfrac53\)，即 \(2a-b-4c=0\). 又由 \(D(\eta)=\dfrac59\)，得 \(E(\eta^{2})=D(\eta)+[E(\eta)]^{2}=\dfrac{10}{3}\)，所以 \(\dfrac{a+4b+9c}{N}=\dfrac{10}{3}\)，即 \(7a-2b-17c=0\). 由 \(b=2a-4c\) 代入第二个方程，得 \(7a-2(2a-4c)-17c=0\)，从而 \(3a-9c=0\)，即 \(a=3c\)，再得 \(b=2c\). 因此 \(a:b:c=3:2:1\).
\end{enumerate}
\end{solution}

\begin{problem}
如图，在四面体 \(A - BCD\) 中，\(AD \perp\) 平面 \(BCD\)，\(BC \perp CD\)，\(AD = 2\)，\(BD = 2\sqrt{2}\). \(M\) 是 \(AD\) 的中点，\(P\) 是 \(BM\) 的中点，点 \(Q\) 在线段 \(AC\) 上，且 \(AQ = 3QC\).

\vspace{-0.5em}

\begin{enumerate}
\item 证明： \(PQ\parallel\) 平面 \(BCD\)；
\item 若二面角 \(C - BM - D\) 的大小为 \(60^{\circ}\)，求 \(\angle BDC\) 的大小.
\end{enumerate}

\bitmapfigure[max width=6.0cm]{img/2013/zhejiang_science/q20_fig1.png}
\end{problem}

\begin{solution}
\begin{enumerate}
\item 建立空间直角坐标系，使平面 \(BCD\) 为 \(z=0\)，取 \(D=(0,0,0)\)，\(A=(0,0,2)\)，\(B=(2\sqrt2,0,0)\)，并设 \(C=(u,v,0)\). 则 \(M=(0,0,1)\)，\(P=(\sqrt2,0,\frac12)\). 由 \(AQ=3QC\)，得 \(Q=\dfrac14A+\dfrac34C=\left(\dfrac{3u}{4},\dfrac{3v}{4},\dfrac12\right)\). 因此 \(P\)，\(Q\) 的竖坐标相同，直线 \(PQ\) 平行于平面 \(z=0\)，即 \(PQ\parallel\) 平面 \(BCD\).

\item 记 \(L=2\sqrt2\)，则 \(B=(L,0,0)\)，\(M=(0,0,1)\). 平面 \(DBM\)，\(CBM\) 的法向量可分别取为
\(\bs{n}_{1}=\overrightarrow{BD}\times\overrightarrow{BM}=(0,L,0)\)，\(\bs{n}_{2}=\overrightarrow{BC}\times\overrightarrow{BM}=(v,L-u,Lv)\). 由二面角为 \(60^{\circ}\)，得
\(\frac{|\bs{n}_{1}\cdot\bs{n}_{2}|}{|\bs{n}_{1}|\,|\bs{n}_{2}|}=\frac{|L-u|}{\sqrt{9v^{2}+(L-u)^{2}}}=\frac12\)，从而 \((L-u)^{2}=3v^{2}\).

又因 \(BC\perp CD\)，有 \((B-C)\cdot(D-C)=0\)，即 \(u^{2}+v^{2}=Lu\)，也就是 \(v^{2}=u(L-u)\). 点 \(C\) 不与 \(B\)，\(D\) 重合，故可除以 \(L-u\)，由 \((L-u)^{2}=3u(L-u)\) 得 \(L-u=3u\). 于是 \(u=\dfrac{L}{4}\)，且 \(v^{2}=3u^{2}\). 所以
\(\cos\angle BDC=\dfrac{\overrightarrow{DB}\cdot\overrightarrow{DC}}{|DB|\,|DC|}=\dfrac{u}{\sqrt{u^{2}+v^{2}}}=\dfrac12\)，故 \(\angle BDC=60^{\circ}\).
\end{enumerate}
\end{solution}

\begin{problem}
如图，点 \(P(0, -1)\) 是椭圆 \(C_1: \frac{x^2}{a^2} + \frac{y^2}{b^2} = 1\)（\(a > b > 0\)）的一个顶点，\(C_1\) 的长轴是圆 \(C_2: x^2 + y^2 = 4\) 的直径，\(l_1\)，\(l_2\) 是过点 \(P\) 且互相垂直的两条直线，其中 \(l_1\) 交圆 \(C_2\) 于 \(A\)，\(B\) 两点，\(l_2\) 交椭圆 \(C_1\) 于另一点 \(D\).

\vspace{-0.5em}

\begin{enumerate}
\item 求椭圆 \(C_1\) 的方程；
\item 求 \(\triangle ABD\) 面积取最大值时直线 \(l_1\) 的方程.
\end{enumerate}

\bitmapfigure[max width=6.0cm]{img/2013/zhejiang_science/q21_fig1.png}
\end{problem}

\begin{solution}
\begin{enumerate}
\item 点 \(P(0,-1)\) 是椭圆的短轴顶点，所以 \(b=1\). 椭圆的长轴长为 \(2a\)，而圆 \(x^{2}+y^{2}=4\) 的直径长为 \(4\)，故 \(2a=4\)，即 \(a=2\). 所以椭圆方程为 \(\dfrac{x^{2}}{4}+y^{2}=1\).

\item 若 \(l_{1}\) 为竖直线，则 \(l_{2}\) 为 \(y=-1\)，它与椭圆只有公共点 \(P\)，不符合题意. 因而设直线 \(l_{1}\) 的斜率为 \(m\)，则 \(l_{1}:y=mx-1\)，与之垂直的直线 \(l_{2}\) 可参数表示为 \((x,y)=(mt,-1-t)\). 直线 \(l_{1}\) 到圆心 \(O\) 的距离为 \(\dfrac{1}{\sqrt{m^{2}+1}}\)，所以
\(|AB|=2\sqrt{4-\dfrac{1}{m^{2}+1}}=2\sqrt{\dfrac{4m^{2}+3}{m^{2}+1}}\).

将 \((x,y)=(mt,-1-t)\) 代入椭圆方程，得 \(\dfrac{m^{2}t^{2}}{4}+(-1-t)^{2}=1\)，即 \(t\left(\left(1+\dfrac{m^{2}}{4}\right)t+2\right)=0\). 其中 \(t=0\) 对应点 \(P\)，另一交点对应 \(t=-\dfrac{8}{m^{2}+4}\)，故 \(|PD|=\dfrac{8\sqrt{m^{2}+1}}{m^{2}+4}\). 因 \(l_{1}\perp l_{2}\)，三角形高为 \(|PD|\)，于是面积
\(S_{\triangle ABD}=\dfrac12|AB|\,|PD|=\dfrac{8\sqrt{4m^{2}+3}}{m^{2}+4}\). 令 \(u=m^{2}\geqslant0\)，则
\(S_{\triangle ABD}^{2}=64\dfrac{4u+3}{(u+4)^{2}}\)，而 \(\left(\dfrac{4u+3}{(u+4)^{2}}\right)'=\dfrac{10-4u}{(u+4)^{3}}\). 因此面积在 \(u=\dfrac52\) 时取得最大值，即 \(m=\pm\sqrt{\dfrac52}=\pm\dfrac{\sqrt{10}}{2}\). 所求直线为 \(l_{1}:y=\pm\dfrac{\sqrt{10}}{2}x-1\).
\end{enumerate}
\end{solution}

\begin{problem}
已知 \(a \in \R\)，函数 \(f(x) = x^3 - 3x^2 + 3ax - 3a + 3\).

\begin{enumerate}
\item 求曲线 \(y = f(x)\) 在点 \((1, f(1))\) 处的切线方程；
\item 当 \(x \in [0, 2]\) 时，求 \(|f(x)|\) 的最大值.
\end{enumerate}
\end{problem}

\begin{solution}
\begin{enumerate}
\item \(f(1)=1\)，且 \(f'(x)=3x^{2}-6x+3a\)，所以 \(f'(1)=3(a-1)\). 曲线在 \((1,f(1))=(1,1)\) 处的切线方程为 \(y-1=3(a-1)(x-1)\)，即 \(y=3(a-1)x-3a+4\).

\item 令 \(t=x-1\)，则 \(-1\leqslant t\leqslant1\)，且
\(f(x)=t^{3}+3(a-1)t+1\)，\(f'(x)=3\bigl(t^{2}+a-1\bigr)\).

当 \(a\leqslant0\) 时，\(t^{2}+a-1\leqslant0\)，函数在 \([-1,1]\) 上单调递减，故最大绝对值为 \(\max\{|f(0)|,|f(2)|\}=\max\{3(1-a),1-3a\}=3(1-a)\).

当 \(0<a<1\) 时，令 \(r=\sqrt{1-a}\)，则 \(0<r<1\)，驻点为 \(t=-r\)，\(t=r\)，并且 \(f(-r)=1+2r^{3}\)，\(f(r)=1-2r^{3}\)，\(f(-1)=3r^{2}\)，\(f(1)=2-3r^{2}\). 其中 \(f(-r)\) 为极大值，\(f(r)\) 为极小值. 先有 \(f(-r)-f(-1)=(1-r)^{2}(2r+1)\geqslant0\). 若 \(f(r)\geqslant0\)，则 \(f(-r)>f(r)=|f(r)|\)；若 \(f(r)<0\)，则 \(f(-r)+f(r)=2>0\)，故 \(f(-r)>|f(r)|\). 又当 \(r>\dfrac12\) 时，若 \(f(1)\geqslant0\)，则 \(f(-r)-f(1)=(2r-1)(r+1)^{2}\geqslant0\)；若 \(f(1)<0\)，则 \(f(-r)+f(1)=3(1-r^{2})+2r^{3}>0\)，所以 \(f(-r)\geqslant|f(1)|\). 另一方面，\(f(1)-f(-r)=(1-2r)(r+1)^{2}\). 因此当 \(0<r\leqslant\dfrac12\)，最大绝对值为 \(f(1)=2-3r^{2}=3a-1\)；当 \(\dfrac12<r<1\)，最大绝对值为 \(f(-r)=1+2r^{3}\).

当 \(a\geqslant1\) 时，\(f'(x)\geqslant0\)，函数在 \([-1,1]\) 上单调递增，最大绝对值为 \(f(1)=3a-1\). 综合上述情形，
\[
\max_{0\leqslant x\leqslant2}|f(x)|=
\begin{cases}
3(1-a),&a\leqslant0,\\
1+2(1-a)^{\frac32},&0<a<\dfrac34,\\
3a-1,&a\geqslant\dfrac34.
\end{cases}
\]
\end{enumerate}
\end{solution}
